\begin{abcliste}
\abc Mit $T_1 = \TEiF = \TEi$ liefert die ideale Gasgleichung
\begin{align}
 n &= \nLuF = \frac{\pEi\cdot\VEi}{\ncR\cdot\TEi} = \nLu \approx \result{\nLuP{-3}{2}} \\
 m &= \mLuF = \nLu\cdot\MLu = \mLug \approx \result{\mLugP{-3}{2}}
\end{align}

\abc Zustand 1: $V_1 = \VEiO$, $p_1 = \pEibP{0}{2}$, $T_1 = \TEi$.

Zustand 2 (isotherm, Boyle-Mariotte):
\begin{align}
 p_2 &= \pZwF = \frac{\pEi\cdot\VEi}{\VZw} = \pZw \approx \result{\pZwbP{0}{2}}
\end{align}
Zustand 3 (isochor bei $V_2 = \VZwO$ auf $T_3 = \TDrF = \TDr$ abgekühlt):
\begin{align}
 p_3 &= \pDrF = \pZw\cdot\frac{\TDr}{\TEi} = \pDr \approx \result{\pDrbP{0}{2}}
\end{align}
Zustand 4 (isotherm bei $T_3$ zurück auf $V_1$):
\begin{align}
 p_4 &= \pViF = \pDr\cdot\frac{\VZw}{\VEi} = \pVi \approx \result{\pVibP{0}{2}}
\end{align}

\begin{center}
\begin{tikzpicture}[>=stealth, xscale=0.05, yscale=1.6]
 \draw[gray!30, xstep=25, ystep=0.5] (0,0) grid (200,4);
 \draw[->] (0,0) -- (215,0) node[right] {$V$ in \si{\centi\meter\cubed}};
 \draw[->] (0,0) -- (0,4.4) node[above] {$p$ in \si{\bar}};
 \foreach \x in {50,100,150,200} \node[below] at (\x,0) {\x};
 \foreach \y in {1,2,3,4} \node[left] at (0,\y) {\y};
 \draw[very thick, red, ->, domain=75:180, samples=60] plot (\x, {270/\x});
 \draw[very thick, ->] (180,1.5) -- (180,0.9485);
 \draw[very thick, blue, <-, domain=75:180, samples=60] plot (\x, {170.73/\x});
 \draw[thick, dashed, ->] (75,2.2763) -- (75,3.6);
 \node[circle, fill, inner sep=1.5pt, label=above right:{1}] at (75,3.6) {};
 \node[circle, fill, inner sep=1.5pt, label=above right:{2}] at (180,1.5) {};
 \node[circle, fill, inner sep=1.5pt, label=below right:{3}] at (180,0.9485) {};
 \node[circle, fill, inner sep=1.5pt, label=below left:{4}] at (75,2.2763) {};
 \node[red, above right] at (120,2.25) {$T_1$};
 \node[blue, below left] at (120,1.42) {$T_3$};
\end{tikzpicture}
\end{center}
Die Zustände 1 und 2 liegen auf der Isotherme $T_1$, die Zustände 3 und 4 auf der Isotherme $T_3$ (Hyperbeln $pV = \text{konst.}$). Das isochore Erwärmen von 4 nach 1 (gestrichelt) schliesst den Kreisprozess des Heissluftmotors (Stirling-Prozess).
\end{abcliste}
